% ============================================================
% 00-mo-dau.tex — Mở đầu (Bình)
% Bốn câu hỏi nghiên cứu nêu đầy đủ ở mục 2; Tiểu kết Chương 1, Chương 3 và Kết luận dùng lại tên ngắn CH1–CH4.
% ============================================================
\chapter*{MỞ ĐẦU}
\addcontentsline{toc}{chapter}{MỞ ĐẦU}

\section*{1. Lý do thực hiện tiểu luận}

Tổ chức tín dụng cần đánh giá khả năng trả nợ của khách hàng để hạn chế nợ xấu. Một bước thường gặp trong đánh giá này là
chia khách hàng thành các nhóm có mức rủi ro khác nhau. Trong khung Basel, ngân hàng dùng cách tiếp cận xếp hạng nội bộ
(IRB) phải xếp các khoản vay vào những nhóm đồng nhất rồi ước lượng tham số rủi ro, và các nghiên cứu như
\cite{bosker2025clustered,ahelegbey2023p2p} cho thấy mô hình dựng riêng cho từng nhóm khách hàng dự báo tốt hơn một mô hình
chung. Khi không dựa vào nhãn vỡ nợ, phân cụm là công cụ tự nhiên để chia nhóm.

Ranh giới giữa các nhóm rủi ro thường không rõ: một khách hàng có thể giống nhóm rủi ro thấp ở một số chỉ tiêu và giống
nhóm rủi ro trung bình ở các chỉ tiêu khác. Phân cụm mờ, với thuật toán tiêu biểu là Fuzzy C-Means (FCM), mô tả được
những khách hàng ``biên'' như vậy bằng độ thuộc vào nhiều cụm thay vì ép họ vào một cụm duy nhất.

FCM có 2 điểm yếu được tài liệu nêu. Điểm yếu được nhắc nhiều nhất là thuật toán chỉ bảo đảm hội tụ tới cực tiểu địa
phương hoặc điểm yên ngựa của hàm mục tiêu~\cite{bezdek1987counterexamples}, nên kết quả có thể phụ thuộc khởi tạo. Phần lớn
các bài ghép FCM với thuật toán tối ưu bầy đàn (PSO) hoặc metaheuristic khác nhắm vào điểm yếu này bằng cách để bầy tìm tâm
cụm~\cite{qian2022psofcm,zhang2025qpsofcm,mane2014hpso}. Tuy vậy Hathaway và Bezdek ghi nhận trong thực nghiệm FCM khá ít
phụ thuộc khởi tạo~\cite{hathaway1988recent} (tr.~242--243), và PSO khi tìm tâm chỉ tối ưu đúng hàm mà FCM vốn tự làm giảm.
Điểm yếu thứ hai ít được nhắc hơn: FCM dùng khoảng cách Euclid nên coi mọi thuộc tính ngang
nhau~\cite{liu2021msgwo}. Hồ sơ tín dụng có nhiều chỉ tiêu đo gần cùng một thứ, mà FCM không tự chọn được trọng số
cho chúng, nên một thuật toán tìm kiếm như PSO có lý do để tham gia.

Nhóm bắt đầu theo cách ghép phổ biến: dùng PSO gốc và 4 phiên bản cải tiến tìm tâm khởi tạo cho FCM trên 2 bộ dữ liệu tín
dụng. Kết quả không như kỳ vọng: mọi phiên bản cho cùng nghiệm với FCM khởi tạo k-means++ mà chậm hơn (tới khoảng 24 lần trên Taiwan Default). Đối chiếu
với tài liệu, các bài đề xuất PSO--FCM thường báo kết quả tốt hơn FCM, nhưng chủ yếu so với FCM khởi tạo ngẫu nhiên và không
cùng ngân sách tính toán; các nghiên cứu so sánh độc lập lại cho thấy PSO--FCM hoà hoặc thua FCM, và chậm hơn
(mục~\ref{sec:tinh-hinh-nghien-cuu}). Phép đo số cực tiểu của hàm mục tiêu giải thích vì sao: trên dữ liệu đã thử, hàm này hầu
như chỉ có 1 cực tiểu, nên PSO tìm tâm không còn gì để tìm. Nhóm vì vậy chuyển hướng: để PSO tối ưu thứ FCM không tự tối ưu
được là trọng số thuộc tính, so 5 phiên bản PSO trong vai trò mới này, và kiểm kết quả trên mẫu độc lập. Tiểu luận giữ
lại hướng ban đầu làm bằng chứng cho lần chuyển hướng.

\section*{2. Mục tiêu và câu hỏi nghiên cứu}

Câu hỏi của tiểu luận là: \emph{\cauHoiTong} Các mục tiêu cụ thể là:
\begin{enumerate}[label=(\alph*), leftmargin=*]
  \item trình bày FCM, thuật toán PSO gốc và 4 phiên bản cải tiến (PSO-W, PSO-S, PSO-LOG, QPSO);
  \item thử hướng phổ biến, dùng PSO tìm tâm khởi tạo cho FCM, trên dữ liệu tín dụng, và đo số cực tiểu của hàm mục tiêu để
  giải thích kết quả;
  \item đề xuất cải tiến: dùng PSO tìm trọng số thuộc tính cho FCM, với fitness không dùng nhãn vỡ nợ, kèm quy trình áp dụng;
  \item thử nghiệm đề xuất trên 7 bộ dữ liệu tín dụng, so với FCM thường, trọng số tính bằng công thức và tìm ngẫu nhiên
  cùng ngân sách, theo khả năng tách nhóm rủi ro, chỉ số phân cụm và thời gian;
  \item so 5 phiên bản PSO trong vai trò tìm trọng số và kiểm trọng số tìm được trên mẫu độc lập.
\end{enumerate}
Câu hỏi này được tách thành 4 câu hỏi nghiên cứu:
\begin{enumerate}[label=\textbf{CH\arabic*.}, leftmargin=*]
  \item \emph{Hướng ban đầu:} \chMot
  \item \emph{PSO tìm trọng số:} \chHai
  \item \emph{So sánh các phiên bản:} \chBa
  \item \emph{Kiểm chứng và giới hạn:} \chBon
\end{enumerate}

\section*{3. Đối tượng và phạm vi nghiên cứu}

Đối tượng nghiên cứu là PSO và các phiên bản cải tiến khi dùng để tìm trọng số thuộc tính cho FCM trong bài toán phân khúc
khách hàng vay, đối chiếu với vai trò tìm tâm khởi tạo. Phạm vi như sau.
\begin{itemize}
  \item Phân cụm \emph{không giám sát}: nhãn vỡ nợ không đi vào thuật toán nào, chỉ dùng để đánh giá phân hoạch sau khi phân
  cụm và, trong quy trình đề xuất, để quyết định có áp trọng số hay không.
  \item Năm thuật toán PSO: PSO gốc (Kennedy và Eberhart, 1995), PSO-W (Shi và Eberhart, 1998, 1999), PSO-S và PSO-LOG (Qian và
  cộng sự, 2022), QPSO (Sun và cộng sự, 2004). Đối chứng của phần trọng số: FCM thường, WFCM với trọng số tính bằng công
  thức~\cite{keller2000weighting}, tìm ngẫu nhiên cùng ngân sách. Đối chứng của phần tìm tâm: 3 cách khởi tạo độ thuộc ngẫu
  nhiên, $c$ điểm dữ liệu ngẫu nhiên và k-means++~\cite{arthur2007kmeanspp}.
  \item Hệ số mờ $m = 2$, số cụm $c$ từ 2 đến 4 cho phần trọng số; phần tìm tâm kiểm tra thêm $m = 1{,}5$ và $c$ đến 8.
  \item Dữ liệu: 2 bộ chính là dữ liệu dự báo tình trạng khó khăn tài chính của doanh nghiệp niêm yết Việt Nam và dữ liệu vỡ nợ thẻ tín
  dụng Taiwan Default; 5 bộ tín dụng công khai khác dùng để kiểm độ rộng của kết quả.
\end{itemize}

\section*{4. Phương pháp nghiên cứu}

Nghiên cứu gồm 2 giai đoạn. Giai đoạn đầu thử PSO tìm tâm khởi tạo theo cách của các bài đã đọc và đo số
cực tiểu của hàm mục tiêu để giải thích kết quả; giai đoạn sau, khi hướng đó không cho kết quả tốt, chuyển sang PSO tìm trọng
số thuộc tính với đầy đủ đối chứng và mẫu kiểm chứng. Hai giai đoạn dùng chung dữ liệu, tiền xử lý, seed và cách kiểm định.
Tiểu luận dùng 2 phương pháp. Thứ nhất là tổng hợp tài liệu: đọc 10 bài chính cùng các bài
nền tảng về FCM, trọng số thuộc tính và PSO. Thứ hai là thực nghiệm có đối chứng. Mọi phương pháp chạy trên cùng dữ liệu đã
tiền xử lý, với cùng 30 seed (hạt giống ngẫu nhiên), cùng tâm khởi tạo và cùng mẫu đánh giá; thời gian thực đo trên cùng máy.
Các phương pháp được so theo cặp seed bằng phép thử Wilcoxon~\cite{wilcoxon1945} kèm hiệu chỉnh Holm~\cite{holm1979}. Trọng
số học trên một mẫu được kiểm trên một mẫu độc lập. Kết quả tái lập được bằng mã nguồn ở Phụ lục~\ref{app:code}.

Về đóng góp, tiểu luận (i) đưa PSO tìm trọng số thuộc tính vào FCM trên dữ liệu tín dụng với 3 lựa chọn thiết kế cho
fitness, (ii) tách phần lợi ích của cơ chế bầy đàn bằng đối chứng tìm ngẫu nhiên cùng ngân sách và đối chứng trọng số giải
tích, (iii) so 5 phiên bản PSO trong vai trò này, và (iv) đề xuất quy trình chọn vai trò cho PSO, bắt đầu bằng chẩn đoán
số cực tiểu của hàm mục tiêu (mục~\ref{sec:de-xuat}).

\section*{5. Bố cục tiểu luận}

Chương~\ref{chap:tong-quan} tổng quan về phân cụm trong đánh giá rủi ro tín dụng, FCM và 2 điểm yếu của nó, PSO và các
phiên bản cải tiến, 2 cách ghép PSO với FCM, và tình hình nghiên cứu. Chương~\ref{chap:phuong-phap} trình bày FCM, PSO gốc và
4 phiên bản cải tiến, các chỉ số đánh giá, rồi theo thứ tự nhóm thực hiện: hướng ban đầu (PSO tìm tâm), trọng số thuộc
tính trong FCM, rồi đề xuất cải tiến của nhóm (PSO tìm trọng số thuộc tính và quy trình áp dụng).
Chương~\ref{chap:thuc-nghiem} mô tả dữ liệu, thiết kế thực nghiệm, trình bày hướng ban đầu và lý do chuyển hướng, rồi trả lời
lần lượt CH1 (mục~\ref{sec:diem-xuat-phat}),
CH2 (mục~\ref{sec:pso-trong-so}), CH3 (mục~\ref{sec:phien-ban-trong-so}) và CH4 (mục~\ref{sec:kiem-chung}), diễn giải các
phân khúc tìm được, và cuối cùng đánh giá chung từng thuật toán qua cả 2 vai trò (mục~\ref{sec:danh-gia-chung}). Phụ lục~\ref{app:code} trình bày mã nguồn kèm cách chạy
lại thực nghiệm. Mục 6 dưới đây ghi phân công của từng thành viên.

\section*{6. Phân công thực hiện}

Nhóm chia việc theo thuật toán: mỗi thành viên phụ trách trọn một mạch gồm bài báo, mục lý thuyết, mã nguồn, thí nghiệm và
mục kết quả tương ứng, để tự trình bày và trả lời câu hỏi về phần của mình (Bảng~\ref{tab:phan-cong}). Phân công đọc hiểu
từng tệp mã nguồn ở Bảng~\ref{tab:pl-phan-cong}.

\begin{table}[htbp]
\centering
\caption{Phân công thực hiện tiểu luận}
\label{tab:phan-cong}
\footnotesize
\setlength{\tabcolsep}{3pt}
\renewcommand{\arraystretch}{1.15}
\begin{tabular}{@{}>{\raggedright\arraybackslash}p{2.2cm}>{\raggedright\arraybackslash}p{3.3cm}>{\raggedright\arraybackslash}p{1.7cm}>{\raggedright\arraybackslash}p{1.5cm}>{\raggedright\arraybackslash}p{2.4cm}>{\raggedright\arraybackslash}p{3.5cm}@{}}
\toprule
Thành viên & Thuật toán & Bài báo & Chương 1 & Chương 2 & Chương 3 \\
\midrule
Vương Quốc Bình (trưởng nhóm) & FCM, khoảng cách có trọng số, WFCM giải tích, quy trình đề xuất & 07, 08 &
\ref{sec:rui-ro-tin-dung}, \ref{sec:phan-cum-mo} & \ref{sec:fcm}, \ref{sec:trong-so-fcm}, \ref{sec:quy-trinh} &
\ref{sec:kiem-chung} (TN5), tiểu kết \\
\addlinespace
Trần Trường Giang & PSO gốc, PSO-W, khung PSO tìm trọng số & 01, 09, 10 & \ref{sec:pso-tong-quan}, \ref{sec:metaheuristic} &
\ref{sec:pso-goc}, \ref{sec:pso-w}, \ref{sec:de-xuat}.1--\ref{sec:de-xuat}.4 & \ref{sec:du-lieu}, \ref{sec:thiet-lap},
\ref{sec:chuyen-huong}, \ref{sec:pso-trong-so} (TN3) \\
\addlinespace
Nguyễn Quang Hà & PSO-S, PSO-LOG, chỉ số đánh giá, kiểm định, đánh giá chung các thuật toán & 02, 03 & \ref{sec:cac-nghien-cuu} & \ref{sec:pso-s-log},
\ref{sec:chi-so} & \ref{sec:kiem-dinh}, \ref{sec:phien-ban-trong-so} (TN4), \ref{sec:dien-giai}, \ref{sec:danh-gia-chung} \\
\addlinespace
Nguyễn Tiến Hải & QPSO, PSO tìm tâm khởi tạo, chẩn đoán số cực tiểu & 04, 05, 06 & \ref{sec:trai-chieu}, \ref{sec:diem-de-ngo} &
\ref{sec:qpso}, \ref{sec:pso-khoi-tao} & \ref{sec:e5-dat-j}, \ref{sec:e5-doi-chung} (TN1), \ref{sec:vi-sao} (TN2) \\
\bottomrule
\end{tabular}
\par\smallskip\parbox{\linewidth}{\scriptsize Phần chung: Bình viết Mở đầu, Kết luận, Tóm tắt và ghép bản LaTeX; Giang soạn
Phụ lục~\ref{app:code}; Hà điều phối hỏi đáp; Hải làm slide thuyết trình.}
\end{table}

